% SPDX-FileCopyrightText: 2026 David Shirvanyants
% SPDX-License-Identifier: CC0-1.0

\chapter{Unsupported-area Detection}
\label{chap:detect}

The Detect command identifies the lowest model surfaces that cannot be reached
from material in the preceding layer. It operates on a snapshot of the
selected models after their transformation matrices have been applied. The
work runs in a background thread, and its result is discarded if the source
models change before completion.

\section{Layer Analysis}

The transformed triangle meshes are combined and sliced on the global build
grid using the configured layer thickness, first-layer offset, segmentation
tolerance, and healing threshold. Each layer is converted to integer-scaled
Clipper2 polygons so that unions, offsets, intersections, and differences use
robust polygon operations. Contour winding preserves solid regions and holes.

The first layer is considered fully supported. For every later layer $i$,
the permissible horizontal reach from the preceding layer is

\[
    r_i = (z_i-z_{i-1})
          \left(c + \frac{1}{\tan\alpha}\right),
\]

where $c$ is the overhang coefficient and $\alpha$ is the critical angle.
The complete material area of layer $i-1$ is expanded by $r_i$. This use of
the complete layer is deliberate: an unsupported region will eventually
receive a support and may therefore carry material in the layers above it.
Consequently, Detect marks the lowest unsupported surface instead of repeating
the same warning through the height of an overhang.

The current layer is separated into connected solid components. A component
is eligible for layer-to-layer support only when it intersects the unexpanded
preceding layer. This test rejects an orphan that appears close to, but is not
actually connected with, an older component. For an eligible component, the
supported portion is its intersection with the expanded preceding material.
The union of these intersections is subtracted from the complete current
layer. The remainder contains hanging edges and orphan components and is the
unsupported area for layer $i$.

\begin{figure}[htbp]
    \centering
    \begin{tikzpicture}[node distance=4mm]
        \node[algstep] (mesh) {Apply model transforms and slice on the global build grid};
        \node[algstep, below=of mesh] (previous) {Expand complete preceding-layer material by $r_i$};
        \node[algstep, below=of previous] (components) {Split the current layer into connected solid components};
        \node[algstep, below=of components] (eligible) {Keep components which intersect the unexpanded preceding layer};
        \node[algstep, below=of eligible] (subtract) {Intersect eligible components with the expanded area, then subtract from the current layer};
        \node[algstep, below=of subtract] (result) {Accumulate unsupported polygons and their absolute area};
        \draw[algarrow] (mesh) -- (previous);
        \draw[algarrow] (previous) -- (components);
        \draw[algarrow] (components) -- (eligible);
        \draw[algarrow] (eligible) -- (subtract);
        \draw[algarrow] (subtract) -- (result);
    \end{tikzpicture}
    \caption{Unsupported-area detection for one layer after slicing.}
\end{figure}

\section{Result}

The output retains the layer positions and unsupported polygons and reports
their total absolute area. The GUI tessellates those polygons for a red-orange
overlay. For visibility, an overlay is drawn at the preceding layer's Z
position; this is only a display offset and does not alter the analysis data.
The same result supplies the score used by orientation optimization and the
input surfaces used by support generation.
